\section{Complex exponential and trigonometric functions}
\label{sec:complexexp}

%mbxINTROSUBSECTION

\sectionnotes{1 lecture}

\subsection{The complex exponential}

Let
\begin{equation*}
E(z) \coloneqq \sum_{k=0}^\infty \frac{1}{k!} z^k .
\end{equation*}
This series converges for all $z \in \C$.
So by
\corref{cor:powerseranalytic}, $E$ is analytic on $\C$.   We notice that $E(0) = 1$,
and that for $z=x \in \R$, $E(x) \in \R$.  Keeping $x$ real, direct
computation shows
\begin{equation*}
\frac{d}{dx} \bigl( E(x) \bigr) = E(x) .
\end{equation*}
In \volIref{\sectionref*{vI-sec:logandexp} of volume I}{\sectionref{sec:logandexp}} (or by Picard's theorem), we proved that
the unique function on the real line satisfying $E' = E$ and
$E(0) = 1$ is the exponential.  In other words, for $x \in \R$, $e^x = E(x)$.

For complex numbers $z$, we define\glsadd{not:complexexp}
\begin{equation*}
e^z \coloneqq E(z) = 
\sum_{k=0}^\infty \frac{1}{k!} z^k .
\end{equation*}
On the real line, this new definition agrees with our previous one.
See \figureref{fig:complexexpgraphs}.  Notice that in the $x$ direction
(the real direction)
the graph behaves like the real exponential, and in the $y$ direction
(the imaginary direction) the graph oscillates.

\begin{myfigureht}
\myincludepdft{real_imag_exp}{%
The surface graphs. 
For a fixed x both graphs are sinusoidal in the y direction
with the amplitude of the oscillations being larger for larger x.
In fact for negative x the graph is so small at this scale that it is
impossible to note the oscillations there.
In the right graph the sinusoidal wave is shifted so that
on the left the graph goes up exponentially above the x-axis
while on the right the graph above the x-axis is simply zero,
that is the x-axis is contained in the graph.  The graph above
x-axis is marked with a bold line on the graph.}
\caption{Graphs of the real part (left) and imaginary part (right)
of the complex exponential $e^z = e^{x+iy}$.  The $x$-axis goes from $-4$ to
$4$, the $y$-axis goes from $-6$ to $6$, and the vertical axis goes from
$-e^{4} \approx -54.6$
to
$e^{4} \approx 54.6$.  The plot of the real exponential ($y=0$)
is marked in a bold line.\label{fig:complexexpgraphs}}
\end{myfigureht}

\begin{prop}[Law of exponents\index{law of exponents}]
Let $z,w \in \C$ be complex numbers.  Then
\begin{equation*}
e^{z+w} = e^z e^w.
\end{equation*}
\end{prop}

\begin{proof}
We already know that
the equality
$e^{x+y} = e^x e^y$ holds for all
real numbers $x$ and $y$.
For every fixed $y \in \R$, consider the expressions as
functions of $x$ and apply the identity theorem
(\thmref{thm:identityanalytic}) to get that
$e^{z+y} = e^ze^y$ for all $z \in \C$.  Fixing an arbitrary $z \in \C$,
we get
$e^{z+y} = e^ze^y$ for all $y \in \R$.  Again by the identity theorem 
$e^{z+w} = e^z e^w$
for all $w \in \C$.
\end{proof}

A simple consequence of the proposition is that $e^z \neq 0$ for all $z \in \C$,
as $e^z e^{-z} = e^{z-z} = 1$.
This computation means that ${(e^z)}^{-1} = e^{-z}$.
Combining that fact with the law of exponents gives
\begin{equation*}
{(e^z)}^{n} = e^{nz} \qquad \text{for all } n \in \Z.
\end{equation*}
A yet more complicated consequence is that we
can compute the power series for the exponential at any point $a \in
\C$:
\begin{equation*}
e^z = e^a e^{z-a} = \sum_{k=0}^\infty \frac{e^a}{k!} {(z-a)}^k .
\end{equation*}

\subsection{Trigonometric functions and $\pi$}

We can now finally define \emph{\myindex{sine}} and \emph{\myindex{cosine}}
by the equation
\begin{equation*}
e^{x+iy} = e^x \bigl( \cos(y) + i \sin(y) \bigr) .
\end{equation*}
In fact, we define sine and cosine for all complex $z$:
\glsadd{not:sin}\glsadd{not:cos}
\begin{equation*}
\cos(z) \coloneqq \frac{e^{iz} + e^{-iz}}{2}
\qquad\text{and}\qquad
\sin(z) \coloneqq \frac{e^{iz} - e^{-iz}}{2i} .
\end{equation*}

Let us use our definition to prove common properties of
sine and cosine.  In the process, we also define the
number $\pi$.

\begin{prop}
The sine and cosine functions have the following properties:
\begin{enumerate}[(i)]
\item For all $z \in \C$,\index{Euler's formula}
\begin{equation*}
e^{iz} = \cos(z) + i\sin(z) \qquad
\text{(Euler's formula)}.
\end{equation*}
\item $\cos(0) = 1$, $\sin(0) = 0$.
\item For all $z \in \C$,
\begin{equation*}
\cos(-z) = \cos(z), \qquad
\sin(-z) = -\sin(z).
\end{equation*}
\item For all $z \in \C$,
\begin{equation*}
\cos(z) = \sum_{k=0}^\infty \frac{{(-1)}^k}{(2k)!} z^{2k} ,
\qquad
\sin(z) = \sum_{k=0}^\infty \frac{{(-1)}^k}{(2k+1)!} z^{2k+1} .
\end{equation*}
\item For all $x \in \R$
\begin{equation*}
\cos(x) = \Re (e^{ix})
\qquad\text{and}\qquad
\sin(x) = \Im (e^{ix}) .
\end{equation*}
\item For all $z \in \C$,
\begin{equation*}
{\bigl( \cos(z) \bigr)}^2 + {\bigl( \sin(z) \bigr)}^2 = 1 .
\end{equation*}
\item For all $x \in \R$,
\begin{equation*}
\babs{\sin(x)} \leq 1, \qquad \babs{\cos(x)} \leq 1 .
\end{equation*}
\item For all $x \in \R$,
\begin{equation*}
\frac{d}{dx} \bigl[ \cos(x) \bigr] = -\sin(x)
\qquad \text{and} \qquad
\frac{d}{dx} \bigl[ \sin(x) \bigr] = \cos(x) .
\end{equation*}
\item For all $x \geq 0$,
\begin{equation*}
\sin(x) \leq x .
\end{equation*}
\item
There exists an $x > 0$ such that $\cos(x) = 0$.  We define
\glsadd{not:pi}
\begin{equation*}
\pi \coloneqq 2 \, \inf \{ x > 0 : \cos(x) = 0 \} .
\end{equation*}
\item
For all $z \in \C$, 
\begin{equation*}
e^{2\pi i} = 1 \qquad \text{and} \qquad e^{z + i 2\pi} = e^z.
\end{equation*}
\item
Sine and cosine are $2\pi$-periodic and not periodic with any smaller
period.  That is, $2\pi$ is the smallest number such that for all $z \in \C$,
\begin{equation*}
\sin(z+2\pi) = \sin(z)
\qquad \text{and} \qquad
\cos(z+2\pi) = \cos(z) .
\end{equation*}
\item
The function $x \mapsto e^{ix}$ is a bijective map from $[0,2\pi)$
onto the set of $z \in \C$ such that $\sabs{z} = 1$.
\end{enumerate}
\end{prop}

The proposition immediately implies that $\sin(x)$ and $\cos(x)$ are 
real whenever $x$ is real.

\begin{proof}
The first three items follow directly from the definition.
The computation of the power series for both is left as an exercise.
As the complex conjugate is a continuous function, the definition
of $e^z$ implies
$\overline{e^z} = e^{\bar{z}}$.  If
$x$ is real,
\begin{equation*}
\overline{e^{ix}} = e^{-ix} .
\end{equation*}
Thus for real $x$,
$\cos(x) =
\frac{e^{ix}+e^{-ix}}{2} =
\frac{e^{ix}+\overline{e^{ix}}}{2} =
\Re (e^{ix})$
and similarly $\sin(x) = \Im (e^{ix})$.

For real $x$, we compute
\begin{equation*}
1 =  e^{ix} e^{-ix}
= e^{ix} \, \overline{e^{ix}}
= \sabs{e^{ix}}^2
= \babs{\cos(x) + i \sin(x)}^2
= {\bigl( \cos(x) \bigr)}^2 + {\bigl( \sin(x) \bigr)}^2 .
\end{equation*}
A slightly more complicated computation shows this fact for
complex numbers, see \exerciseref{exercise:cossinidentity}.
In particular, $e^{ix}$ is \emph{\myindex{unimodular}} for real $x$;
the values lie on the unit circle.
A square of a real number is always nonnegative:
\begin{equation*}
{\bigl(\sin(x)\bigr)}^2 = 1-{\bigl(\cos(x)\bigr)}^2 \leq 1 .
\end{equation*}
So $\babs{\sin(x)} \leq 1$ and similarly 
$\babs{\cos(x)} \leq 1$.

We leave the computation of the derivatives to the reader as exercises.
Let us prove that $\sin(x) \leq x$ for $x \geq 0$.
Consider
$f(x) \coloneqq x-\sin(x)$ and differentiate:
\begin{equation*}
f'(x) = \frac{d}{dx} \bigl[ x - \sin(x) \bigr]
=
1 -\cos(x) \geq 0 ,
\end{equation*}
for all $x \in \R$ as $\babs{\cos(x)} \leq 1$.
In other words, $f$ is increasing and $f(0) = 0$.
So $f$ must be nonnegative when $x \geq 0$ and hence, $\sin(x) \leq x$.

Next, we claim there exists a positive $x$ such that $\cos(x) = 0$.
As $\cos(0) = 1 > 0$, $\cos(x) > 0$
for $x$ near $0$.  Namely,
there is some
$y > 0$ such that $\cos(x) > 0$ on $[0,y)$.
Then $\sin(x)$ is strictly
increasing on $[0,y)$.  As $\sin(0) = 0$, we have
$\sin(x) > 0$ for $x \in (0,y)$.  Take $a \in (0,y)$.  By
the mean value theorem, there is a $c \in (a,y)$ such that
\begin{equation*}
2 \geq \cos(a)-\cos(y) = \sin(c)(y-a) \geq \sin(a)(y-a) .
\end{equation*}
As $a \in (0,y)$, we have $\sin(a) > 0$ and so
\begin{equation*}
y \leq \frac{2}{\sin(a)} + a .
\end{equation*}
Hence there is some largest $y$ such that $\cos(x) > 0$ in $[0,y)$,
and let $y$ be the largest such number.
By continuity, $\cos(y) = 0$.
In fact, $y$ is the
smallest positive $y$ such that $\cos(y) = 0$.  As mentioned,
$\pi$ is defined to be $2y$.

As $\cos(\nicefrac{\pi}{2}) = 0$, we find
${\bigl(\sin(\nicefrac{\pi}{2})\bigr)}^2 = 1$.
As $\sin$ is positive on $(0,\nicefrac{\pi}{2})$, we have
$\sin(\nicefrac{\pi}{2}) = 1$.
Hence,
\begin{equation*}
e^{i \pi /2} = i ,
\end{equation*}
and by the law of exponents,
\begin{equation*}
e^{i \pi} = -1 ,
\qquad 
e^{i 2\pi} = 1 .
\end{equation*}
So $e^{i2\pi} = 1 = e^0$.  The law of exponents also says
\begin{equation*}
e^{z+i2\pi} = e^z e^{i2\pi} = e^z
\end{equation*}
for all $z \in \C$.  Immediately, we also obtain
$\cos(z+2\pi) = \cos(z)$ and $\sin(z+2\pi) = \sin(z)$.
So $\sin$ and $\cos$ are $2\pi$-periodic.

We claim that $\sin$ and $\cos$ are not periodic with a smaller period.  It
suffices to show that if $e^{ix} = 1$ for the
smallest positive $x$, then
$x = 2\pi$.  Let $x$ be the smallest positive $x$ such that
$e^{ix} = 1$.
Of course, $x \leq 2\pi$.
By the law of exponents,
\begin{equation*}
{\bigl(e^{ix/4}\bigr)}^4 = 1 .
\end{equation*}
If $e^{ix/4} = a+ib$, then
\begin{equation*}
{(a+ib)}^4
=a^4-6a^2b^2+b^4 + i\bigl(4ab(a^2-b^2)\bigr)
=1 .
\end{equation*}
Then either $a = 0$ or $a^2 = b^2$.
As $\nicefrac{x}{4} \leq \nicefrac{\pi}{2}$, we have $a = \cos(\nicefrac{x}{4}) \geq 0$ and
$b = \sin(\nicefrac{x}{4}) > 0$.
If $a^2=b^2$, then
$a^4-6a^2b^2+b^4 = -4a^4 < 0$ and in particular not equal to 1.
Therefore $a=0$, in which case $\nicefrac{x}{4} = \nicefrac{\pi}{2}$.
Hence $2\pi$ is the smallest period we could choose for $e^{ix}$
and so also for $\cos$ and $\sin$.

Finally, we wish to show that $e^{ix}$ is one-to-one and onto
from the set $[0,2\pi)$ to the set of $z \in \C$ such that
$\sabs{z} = 1$.  Suppose $e^{ix} = e^{iy}$ and 
$x > y$.  Then
$e^{i(x-y)} = 1$, meaning $x-y$ is a multiple of $2\pi$ and hence
only one of them can live in $[0,2\pi)$.
To show $e^{ix}$ is onto, pick $(a,b) \in \R^2$ such that $a^2+b^2 = 1$.
Suppose first that $a,b \geq 0$.  By the intermediate value theorem,
there must exist an $x \in [0,\nicefrac{\pi}{2}]$ such that
$\cos(x) = a$, and hence $b^2 = \bigl(\sin(x)\bigr)^2$.  As
$b$ and $\sin(x)$ are nonnegative, $b = \sin(x)$.
Since $-\sin(x)$ is the derivative of $\cos(x)$
and $\cos(-x) = \cos(x)$,
we have that $\sin(x) < 0$ for $x \in [\nicefrac{-\pi}{2},0)$.
Using the same reasoning, we obtain that
if $a > 0$ and $b \leq 0$, we can find an $x$ in $[\nicefrac{-\pi}{2},0)$,
and by periodicity,
$x \in [\nicefrac{3\pi}{2},2\pi)$ such that $\cos(x) = a$ and $\sin(x)=b$.
Multiplying by $-1$ is the same as multiplying by $e^{i\pi}$ or
$e^{-i\pi}$.  So we can always assume that $a \geq 0$ (details are left
as an exercise).
\end{proof}

\subsection{The unit circle and polar coordinates}

The arclength of a curve parametrized by $\gamma \colon [a,b] \to \C$ is given
by
\begin{equation*}
\int_a^b \babs{\gamma'(t)} \, dt .
\end{equation*}
We have that $e^{it}$ parametrizes the circle for $t$ in $[0,2\pi)$.
As $\frac{d}{dt} \bigl( e^{it} \bigr) = ie^{it}$, the circumference of the
circle (the arclength) is
\begin{equation*}
\int_0^{2\pi} \sabs{i e^{it}}  \,  dt
=
\int_0^{2\pi} 1  \,  dt  = 2\pi .
\end{equation*}

More generally, $e^{it}$ parametrizes the circle by arclength.
That is, $t$ measures the arclength on a circle of radius 1 by
the angle in radians.  So the definitions of $\sin$ and $\cos$ given
above agree with the standard geometric definitions.

All the points on the unit circle can be achieved by
$e^{it}$ for some $t$.
Therefore,
we can write
a complex number $z \in \C$
(in what we call \emph{\myindex{polar coordinates}}) as
\begin{equation*}
z = r e^{i\theta}
\end{equation*}
for some $r \geq 0$ and $\theta \in \R$.  The $\theta$ is, of course,
not unique as $\theta$ or $\theta+2\pi$ gives the same number.
The law of exponents $e^{a+b} = e^a e^b$ leads to a useful formula for powers
and products of complex numbers in polar coordinates:
\begin{equation*}
{(r e^{i\theta})}^n
= r^n e^{i n \theta} ,
\qquad
(r e^{i\theta})
(s e^{i\gamma})
=
rs e^{i(\theta+\gamma)} .
\end{equation*}

